% Producer of the pairwise receipt (G3) of subsec:trace-rounded-comparison.
% Issue #372 (parent #56, C15). Correspondence and evidence:
% blueprint/units/YM-000026/handoff/results/C15/author/hairer-claude-372-20261114/.
\subsection{The pairwise receipt from the level budget}
\label{subsec:trace-rounded-receipt}

We now derive the receipt (G3) of
Subsection~\ref{subsec:trace-rounded-comparison} from two properties of
the funded chart witnesses of
Section~\ref{sec:ex-lattice-firesky-gallery}: a scale allowance, stated in
(A) below, and a root-test bound at each level whose size is an explicit
negative power of the clock monitor, stated in (B) below. The width clause
uses only (A) and the growth of the monitor. The two root-test bounds
\eqref{eq:trace-rounded-receipt-one}--\eqref{eq:trace-rounded-receipt-two}
use only (B) and the growth of the monitor. The selection of the chart
levels then supplies (G3) at the cost $\varepsilon_i=(i+1)^{-1}$. At the
end of the subsection we show how the replica-root estimate of
Section~\ref{sec:rank-root-normalization} yields (B), and list what
remains open.

\paragraph{Charts at a fixed level.}
Throughout, $\mu_k\ge1$ is the clock monitor of (G1), with
$\mu_k\to\infty$; it is the diverging monitor $\mathcal M_k$ of
Section~\ref{sec:rank-root-normalization}. Fix a grade cutoff $d$ and
let
\[
 M=M_d=d+2
\]
be the marker allowance of the funded chart family at grade $d$: the
maximal number of markers per cell in the replica estimate of
Section~\ref{sec:rank-root-normalization}. (It is not the arity cap $l$
of Section~\ref{sec:ex-lattice-firesky-gallery}.) Consider a chart of the
family at level $k$ with base arc $b>0$. Put
\[
 m=\mu_k^{\,b},\qquad r=a_km,
\]
so that $r=g_b(k)$ in the notation of \eqref{eq:ex-lattice-firesky-gap}.
Let $a_{\ell}$ be the spacing of the middle cell of the chart. Let $w$ be
the root depth of the chart: the depth of the centered tiling on which
the root tests of its comparison receipt are taken, and at which its
chamber condition is asserted. A configuration $x\in(\mathbb R^4)^2$ is
\emph{fit for the chart} if the grid configuration $a_k\lceil mx\rceil$
is coordinatewise separated and each of its coordinates has absolute
value at most $a_\ell/8$. For the chart selected at index $i$, with
$(k,b)=(k_i,b_i)$, this is exactly fitness at level $i$ as defined in
Subsection~\ref{subsec:trace-rounded-comparison}. As there, for a fit
$x$ and $e,e'\in E_{d'}$ with $d'\le d$, we write $f_0^{N,p}$,
$f_1^{N,p}$ for the finite-volume insertions of the chart labels of $e$
and $e'$ at the two sites of $a_k\lceil mx\rceil$. We write
$j_0,j_1,\kappa$ for the reference values at the rounded configuration
$\rho_m(x)$, and $\Phi_s$ ($s\subseteq\{0,1\}$) for the parts of the
corrected polynomial defined there. Thus $\Phi_{\{c\}}=f_c^{N,p}-j_c$ and
$\Phi_{\{0,1\}}=(f_0^{N,p}-j_0)(f_1^{N,p}-j_1)-\kappa$.

We use the following two properties of the funded chart witnesses. In
both, the level threshold is fixed before the chart is chosen.
\begin{enumerate}
\item[(A)] \emph{Allowance.} There is $K_A=K_A(d)$ such that every chart
 of the family at a level $k\ge K_A$ with base arc $b$ satisfies
 \begin{equation}\label{eq:trace-receipt-allowance}
  m\,\mu_k^{2}\le\frac{a_\ell}{a_k},
  \qquad\text{equivalently}\qquad r\,\mu_k^2\le a_\ell .
 \end{equation}
\item[(B)] \emph{Level budget.} There is $K_B=K_B(d)$ with the following
 property. Let $k\ge K_B$ and consider any chart of the family at level
 $k$, any grade $d'\le d$, any basis vectors $e,e'$ of $E_{d'}$, any
 $N>k$, any period $p$, and any $x$ fit for the chart with
 $\rho_m(x)\in\mathcal W$. Then, for every nonempty
 $s\subseteq\{0,1\}$, the root test of the centered tiling at the root
 depth $w$ of the chart, for the Wilson law at $(N,p)$, satisfies
 \begin{equation}\label{eq:trace-receipt-budget}
  \|\Phi_s\|_{\mathrm{root}}\le\mu_k^{-1/(16M)} .
 \end{equation}
\end{enumerate}
The thresholds $K_A$ and $K_B$ may depend on the grade cutoff, the chosen
bases and the window $\mathcal W$. They do not depend on the chart, on
$(N,p)$, or on $x$. Since there are finitely many grades $d'\le d$ and
finitely many pairs of basis vectors, a threshold for each fixed pair
yields one for all of them by taking the maximum. This chart-uniform form
is used in the proof of Proposition~\ref{prop:trace-rounded-receipt},
where the threshold must be known before the producer is called. A
weaker form, with a threshold allowed to depend on the level curve
returned by the producer, would suffice for the first assertion of that
proposition, but not for the argument given for its second assertion.

\begin{lemma}[The width clause]\label{lem:trace-receipt-width}
Let $R>0$. Consider a chart at a level $k\ge K_A$ with base arc $b$, and
suppose $\mu_k^2\ge8R$. Let $t\in(\mathbb Z^4)^2$ be such that
$y=m^{-1}t$ is coordinatewise separated and $|y|\le R$. Then $y$ is fit
for the chart. In particular, if $\mu_k^2\ge800$, the width clause of
(G3) holds for that chart.
\end{lemma}
\begin{proof}
Since $my=t$ has integer coordinates, $\lceil my\rceil=t$, and the grid
configuration of $y$ is
\[
 z=a_k\lceil my\rceil=a_kt=a_kmy=ry .
\]
Multiplication by $r>0$ preserves the strict order of real numbers, so
$z$ is coordinatewise separated because $y$ is. Each coordinate
satisfies $|z_{c,\nu}|=r|y_{c,\nu}|\le r|y|\le rR$, because a coordinate
of a vector in $\mathbb R^{8}$ is bounded by its Euclidean norm. By
\eqref{eq:trace-receipt-allowance}, which holds because $k\ge K_A$, and
by $\mu_k^2\ge8R$,
\[
 rR\le\frac{a_\ell}{\mu_k^2}\,R\le\frac{a_\ell}{8R}\,R=\frac{a_\ell}{8}.
\]
Hence $y$ is fit. The width clause of (G3) is the case $R=100$, and
there $8R=800$.
\end{proof}

\begin{lemma}[Cost of the root tests]\label{lem:trace-receipt-cost}
Assume (B), and let $0<\varepsilon\le1$. If $k\ge K_B$ and
$\mu_k\ge\varepsilon^{-16M}$, then for every chart of the family at level
$k$ the bounds
\eqref{eq:trace-rounded-receipt-one}--\eqref{eq:trace-rounded-receipt-two},
with the root test taken at the root depth $w$ of the chart, hold with
$\varepsilon$ in place of $\varepsilon_i$. This holds for every
grade $d'\le d$, all basis vectors $e,e'$ of $E_{d'}$, all $N>k$, all
$p$, and every $x$ fit for the chart with $\rho_m(x)\in\mathcal W$.
\end{lemma}
\begin{proof}
The function $t\mapsto t^{-1/(16M)}$ is decreasing on $(0,\infty)$, so
$\mu_k\ge\varepsilon^{-16M}$ gives
$\mu_k^{-1/(16M)}\le(\varepsilon^{-16M})^{-1/(16M)}=\varepsilon$. By
(B), $\|\Phi_s\|_{\mathrm{root}}\le\varepsilon$ for each nonempty
$s\subseteq\{0,1\}$. By the identities recalled above, the left sides
of \eqref{eq:trace-rounded-receipt-one} and
\eqref{eq:trace-rounded-receipt-two} are exactly
$\|\Phi_{\{c\}}\|_{\mathrm{root}}$ and
$\|\Phi_{\{0,1\}}\|_{\mathrm{root}}$.
\end{proof}

Since $\mu_k\to\infty$, the conditions
\[
 k\ge K_A(d),\qquad k\ge K_B(d),\qquad \mu_k^2\ge800,\qquad
 \mu_k\ge\varepsilon^{-16M_d}
\]
hold for all sufficiently large $k$. We write $K(d,\varepsilon)$ for the
least $K$ such that they hold for every $k\ge K$. (The monitor is not
assumed monotone, so this is the right form of the threshold.)

\begin{proposition}[Producer of the pairwise receipt]
\label{prop:trace-rounded-receipt}
Assume (A) and (B) at every grade cutoff. Identify the root depth $w_i$
of (G4) with the root depth $w$ of the chart selected at index $i$; that
is, (G4) is required at that depth. Suppose that for every index $i$ the
gallery selection of Section~\ref{sec:ex-lattice-firesky-gallery}, with
grade cutoff $d=i$ and cost index $l=i$, takes its chart at a level
\begin{equation}\label{eq:trace-receipt-level}
 k_i\ge K\bigl(i,(i+1)^{-1}\bigr).
\end{equation}
Then the selected data satisfy (G3) at the cost
$\varepsilon_i=(i+1)^{-1}$, including its width clause. Moreover
the selection can always be made so that \eqref{eq:trace-receipt-level}
holds for every $i$, without changing the conclusions of
Lemma~\ref{lem:ex-lattice-firesky-selection}.
\end{proposition}
\begin{proof}
Fix $i$, and let $M=M_i=i+2$. The selected chart has level
$k_i\ge K(i,(i+1)^{-1})$ and base arc $b_i$, so
$m=\mu_{k_i}^{b_i}=m_i$ and the fitness of the previous paragraphs is
fitness at level $i$. By the identification in the statement, the root
test of (G3) at index $i$ is taken at the depth $w_i=w$ of the selected
chart. Since $k_i\ge K_A(i)$ and $\mu_{k_i}^2\ge800$,
Lemma~\ref{lem:trace-receipt-width} with $R=100$ gives the width clause:
every $y=m_i^{-1}t$ that is coordinatewise separated with $|y|\le100$ is
fit at level $i$. Since $k_i\ge K_B(i)$ and
$\mu_{k_i}\ge(i+1)^{16M}=\varepsilon_i^{-16M}$,
Lemma~\ref{lem:trace-receipt-cost} with $\varepsilon=\varepsilon_i$ gives
\eqref{eq:trace-rounded-receipt-one}--\eqref{eq:trace-rounded-receipt-two}
at depth $w_i$. These hold for every grade $d'\le i$, all basis vectors
of $E_{d'}$, all $N>k_i$, all $p$, and every $x$ fit at level $i$ with
$\rho_{m_i}(x)\in\mathcal W$. As $i$ was arbitrary, this is (G3).

For the second assertion we use only the output of the one-step producer
as stated in Section~\ref{sec:ex-lattice-firesky-gallery}. The chart it
returns at level $k$ and base arc $b>5$ has exact radius equal to the
clock gap $g_b(k)=a_k\mu_k^b$ of \eqref{eq:ex-lattice-firesky-gap}, and the
producer accepts every threshold $\theta>0$ for that radius. Fix $i$ and
put $K=K(i,(i+1)^{-1})$; by (A) and (B) as stated, $K$ is determined
before the producer is called. Since $\mu_k\ge1$ and $b>0$, the radius
satisfies $g_b(k)\ge a_k$. Given the threshold $\theta>0$ required by the
selection at step $i$, call the producer instead with
\[
 \theta'=\min\Bigl(\theta,\ \min_{j<K}a_j\Bigr)>0
\]
(with $\theta'=\theta$ if $K=0$); this minimum is positive because
there are finitely many $j<K$ and each $a_j>0$. If the returned level $k$
satisfied $k<K$, the returned radius would satisfy
$g_b(k)\ge a_k\ge\theta'$, contradicting $g_b(k)<\theta'$. Hence
$k\ge K$, which is \eqref{eq:trace-receipt-level}. Since
$\theta'\le\theta$, the returned output also meets the threshold $\theta$
required in the proof of Lemma~\ref{lem:ex-lattice-firesky-selection},
and every other property of the output is unchanged. Thus that lemma
applies verbatim with the reduced thresholds, and its conclusions hold
for the resulting selection.
\end{proof}

\paragraph{Supply of the level budget.}
The bound \eqref{eq:trace-receipt-budget} is the replica-root estimate of
Section~\ref{sec:rank-root-normalization} at marker allowance $M$. Write
$q_0$ for the number of replica cells, which is also the site count of
the selected output torus. With
\[
 I_k=\frac{\log\mu_k}{16M},\qquad a=e^{2I_k}=\mu_k^{1/(8M)}\ge1 ,
\]
a coefficient whose pushed marked density has the local estimate with
coefficient price $a^{q_0|I|}$ and per-site expense $X\le1+M$ satisfies
\eqref{eq:rank-root}, whose right side is
$\exp\{P_R(X)-4D\log g/q_0\}/a^{|I|}$. The harmonic term
$-4D\log g/q_0$ is nonnegative and must be paid for. That section shows
\[
 P_R(X)-4D\log g/q_0\le T_M
\]
under the following conditions, all at the output level $m$ of the clock:
the coupling bound $\beta\le4\widehat\beta_m$; the harmonic condition
$8D\widehat\beta_m/q_0\le1$, which the clock supplies eventually; the
same clip and row bound $A$ as the clock stage, with
$\theta_R\le\theta_{\rm clock}$, $\Lambda\le\Lambda_{\rm clock}$ and
$\delta\le1$; and the late-scale purchases
$\widehat\beta_m\theta_{\rm clock}^3\le1$,
$\widehat\beta_m\theta_{\rm clock}^4\le1$ and
$\Lambda_{\rm clock}\theta_{\rm clock}^2\le1$. If moreover the stage ticket
satisfies $T_M\le I_k$, then the right side of \eqref{eq:rank-root} is
at most $e^{I_k}/a^{|I|}$. For a nonempty marker subset $|I|\ge1$ and
$a\ge1$, so this is at most $e^{I_k}/a=e^{-I_k}=\mu_k^{-1/(16M)}$. The
marker subsets $I$ are the nonempty $s\subseteq\{0,1\}$; since
$M\ge2$, they are admissible. The coefficients are the one-point and
corrected two-point parts $\Phi_s$ themselves, as
Section~\ref{sec:rank-root-normalization} requires.

The coefficient price $a^{q_0|I|}$ with $X\le1+M$ on the same pushed
density is obtained in Section~\ref{sec:allarity-marked-propagation},
from the scalar envelope \eqref{eq:allarity-density-envelope} at the
constant marker weight \eqref{eq:allarity-marker-weight} with $u=M$.
That derivation assumes the following, each on the common history of the
chart witness:
\begin{itemize}
\item the one-step purchases
 \eqref{eq:allarity-step-purchases}--\eqref{eq:allarity-step-small};
\item the quantitative outputs of the one-step marked compiler at each
 level: remainder price at most $(\alpha+1)z$, principal incidence price
 at most $H^{-2}$, forced price at most $G_i^2$, and discrepancy at most
 $y$ from the next prescribed connected polynomial on every required
 healthy region;
\item the terminal-source purchase that gives the initial representation
 regular price $d_0$ and forced price at most $G_k$;
\item the step-count purchase $(\alpha+2)^{k-l}\le H$;
\item completed propagation to the selected output torus with exactly
 $q_0$ sites, with $L,\ell\le1$ there.
\end{itemize}

\paragraph{Status.}
Proposition~\ref{prop:trace-rounded-receipt} produces (G3) from (A) and
(B), with $w_i$ the root depth of the selected chart. The steps from
(A), (B) and the divergence of the monitor to (G3) are proved above,
including the arrangement of the selected levels. The arithmetic from the
displayed premise $P_R(X)-4D\log g/q_0\le T_M$, the purchase
$T_M\le I_k$ and the coefficient price $a^{q_0|I|}$ with $X\le1+M$ to
\eqref{eq:trace-receipt-budget} is the computation just given. The
following inputs remain open.
\begin{enumerate}
\item \emph{The allowance} (A). It is a scale constraint of the chart at
 its level. It belongs to the funded chart witness of
 Section~\ref{sec:ex-lattice-firesky-gallery}, but that section neither
 displays this inequality nor defines the middle cell whose spacing
 $a_\ell$ enters it. It must be exported from the witness, with its
 threshold $K_A(d)$.
\item \emph{The level budget} (B), at the root depth $w$ of each chart,
 for every $k\ge K_B(d)$ and every chart witness of the funded family at
 level $k$. By the supply paragraph, this requires the pushed marked
 density of the two-site corrected polynomial on the centered tiling at
 depth $w$ to satisfy three conditions:
 \begin{enumerate}
 \item[(a)] it has the coefficient price $a^{q_0|I|}$ with
  $a=\mu_k^{1/(8M)}$ and $X\le1+M$, on a lawful representation in the
  sense of Section~\ref{sec:rank-root-normalization};
 \item[(b)] the clock-side conditions listed in the supply paragraph hold
  at the output level, including the eventual harmonic condition, with
  that eventual threshold expressed as a level threshold in $k$ and
  absorbed into $K_B(d)$;
 \item[(c)] the stage ticket satisfies $T_M\le I_k$.
 \end{enumerate}
 All three must hold uniformly in $N>k$, in $p$, in the fit $x$ with
 rounding in $\mathcal W$, and in the basis labels of grades at most $d$.
\item \emph{Root depth.} The chamber condition (G4) must hold at the same
 root depth $w$ at which (B) is proved, as required by the
 identification in Proposition~\ref{prop:trace-rounded-receipt}.
\end{enumerate}
For (a), the first unsupported implication is from the finite-volume
chart witness (its chart, invoice, legacy-cost, actuation and stamp
data) to the inputs of Section~\ref{sec:allarity-marked-propagation}
listed above, on the history of that witness. We have not located a
derivation of this implication in the present text.
